\section{Smooth Manifolds}

\subsection{Topological manifolds}

\bdefn

    A {\emphcolor topological manifold of dimension $n$} (or an {\emphcolor $n$-manifold}) is a topological space $M$ such that
    \benum
        \item $M$ is Hausdorff (every two points have disjoint open neighborhoods),
        \item $M$ is second-countable (there is a countable basis for $M$),
        \item $M$ is {\emphcolor locally Euclidean} of dimension $n$: every point is contained in a neighborhood homeomorphic to an open subset of $\bR^n$.
    \eenum

\edefn

The empty set is a topological manifold of every dimension, but for the most part we will ignore this special case.
It can be shown that the dimension of a topological manifold is well-defined: its dimension is unique (other than the empty space).
This can be shown via standard tools in algebraic geometry.

The most basic example of an $n$-manifold is $\bR^n$ itself: it is clearly locally Euclidean and Hausdorff, and the set of all balls with rational center and radius forms
a countable basis.

\bdefn

    Let $M$ be an $n$-manifold.
    A {\emphcolor coordinate chart} on $M$ is a pair $(U,\phi)$ with $U\subseteq M$ and $\phi\colon U\to\bR^n$ an embedding (equivalently $\phi\colon U\to\phi(U)\subseteq\bR^n$
    a homeomorphism).
    $U$ is the {\emphcolor coordinate domain} or {\emphcolor coordinate neighborhood} of the chart, $\phi$ is the {\emphcolor (local) coordinate map} the coordinates of $\phi$
    (i.e.\ the composition with the projections $\bR^n\to\bR$) are called the {\emphcolor local coordinates} on $U$.

    By the definition of a manifold, every point in $M$ is contained in the domain of some chart.
    A chart $(U,\phi)$ is {\emphcolor centered} at $p\in U$ if $\phi(u)=0$.
    Clearly every point in $M$ is the center of some chart.

\edefn

\bexam

    Let $U\subseteq\bR^n$ be open and $f\colon U\to\bR^k$ continuous.
    Define its {\emphcolor graph}:
    $$ \Gamma(f) = \set{(x,y)\in\bR^n\times\bR^k}[x\in U,\ f(x)=y], $$
    and endow it with the subspace topology of $\bR^n\times\bR^k$.

    Now consider the projection $\pi_1\colon\bR^{n+k}\to\bR^n$ and its restriction to $\Gamma(f)$: $\phi\colon\Gamma(f)\to U$.
    This is a continuous map, and it has a continuous inverse $\phi^{-1}(x)=(x,f(x))$.
    Thus $\phi$ forms a homeomorphism between $\Gamma(f)$ and $U$, meaning $\Gamma(f)$ is an $n$-manifold.

\eexam

\bexam

    For an integer $n\geq0$, consider the unit $n$-sphere $S^n$ of all vectors $x\in\bR^{n+1}$ with unit norm, endowed with the subspace topology.
    As a subspace of $\bR^{n+1}$, it is Hausdorff and second-countable.
    It is also locally Euclidean of dimension $n$, and thus an $n$-manifold.

    Indeed, define for each $i=1,\dots,n+1$:
    $$ U_i^+ = \set{(x^1,\dots,x^{n+1})\in\bR^{n+1}}[x^i>0],\qquad U_i^- = \set{(x^1,\dots,x^{n+1})\in\bR^{n+1}}[x^i<0] $$
    And consider $f\colon B^n\to\bR$ by $f(u)=\sqrt{1-\norm u^2}$, where $B^n$ is the unit ball on $\bR^n$ (all vectors with norm $<1$).
    Then $U_i^+\cap S^n$ is the graph of the function
    $$ x^i = f(x^1,\dots,\hat x^i,\dots,x^{n+1}), $$
    where the hat means the indicated index is removed, viewed as a function $\bR^n\to\bR$.
    Similarly $U_i^-\cap S^n$ is the graph of $x^i=-f(x^1,\dots,\hat x^i,\dots,x^{n+1})$.

    Thus each $U_i^\pm\cap S^n$ is locally Euclidean of dimension $n$, as per our previous example.
    The maps $\phi_i^\pm\colon U_i^\pm\cap S^n\to B^n$ from the example are given by
    $$ \phi_i^\pm(x^1,\dots,x^{n+1}) = (x^1,\dots,\hat x^i,\dots,x^{n+1}) . $$
    These are all local homeomorphisms to an open subset of $\bR^n$; and furthermore every point in $S^n$ is in some $U_i^\pm\cap S^n$, thus $S^n$ is indeed locally Euclidean
    of dimension $n$.

\eexam

\bexam

    The {\emphcolor $n$-dimensional real projective space}, denoted $\RP^n$ is the quotient of $\bR^{n+1}-0$ by equating scalar multiples of vectors.
    Equivalently, it is the set of $1$-dimensional subspaces of $\bR^{n+1}$.
    Its quotient topology is given by the natural map $\pi\colon\bR^{n+1}-0\to\RP^n$.
    Recall that this means the open sets are those $U\subseteq\RP^n$ for which $\pi^{-1}(U)\subseteq\bR^{n+1}-0$ is open (equivalently if $\pi^{-1}(U)\cup0\subseteq\bR^{n+1}$ is open).

    $\RP^n$ is an $n$-manifold.
    It can be shown to be second-countable (using any countable basis of $\bR^{n+1}$) and Hausdorff.
    We show now local Euclidean-ness.

    For $i=1,\dots,n+1$ define $\tilde U_i\subseteq\bR^{n+1}$ to be the set of all vectors whose $i$th component is nonzero, and define $U_i=\pi(\tilde U_i)$.
    Notice that $\pi^{-1}(U_i)=\tilde U_i$, which is open, thus $U_i\subseteq\RP^n$ is open and $\pi|_{\tilde U_i}\colon\tilde U_i\to U_i$ is a quotient map.
    Now define $\phi_i\colon U_i\to\bR^n$ by
    $$ \phi_i[\bar x] = \frac1{x^i}(x^1,\dots,\hat x^i,\dots,x^n) , $$
    where $[\bar x]$ is the equivalence class of $\bar x\in\bR^{n+1}$.
    $\phi_i$ is continuous since $\phi_i\circ\pi$ is, and it is a quotient map.

    $\phi_i$ has an inverse:
    $$ \phi_i^{-1}(\bar u) = [u^1,\dots,u^{i-1},1,u^i,\dots,u^n], $$
    and this map is continuous so $\phi_i\colon U_i\to\bR^n$ is a homeomorphism.
    And as before, $U_1,\dots,U_{n+1}$ cover $\RP^n$ so it is locally Euclidean of dimension $n$.

\eexam

\bprop

    The product of manifolds is a manifold.
    More specifically if $M_1,\dots,M_k$ are manifolds of dimension $n_1,\dots,n_k$ respectively, then their product space $M_1\times\cdots\times M_k$ is a manifold
    of dimension $n_1+\cdots+n_k$.

\eprop

\bproof

    It is clearly Hausdorff and second countable.
    Let $\bar p\in\prod_1^kM_i$ and let $(U_i,\phi_i)$ be a coordinate chart for $p_i\in M_i$.
    These assimilate into a coordinate chart $(\prod_1^kU_i,\prod_1^k\phi_i\colon\prod_1^kU_i\to\bR^{n_1+\cdots+n_k})$.
    \qed

\eproof

Recall that the {\emphcolor $n$-torus} is the $n$-fold product $T^n=S^1\times\cdots\times S^1$, and by this proposition and the above example is an $n$-manifold.

We now turn to some topological properties of topological manifolds.
Recall the following definitions:

\benum

    A topological space $X$ is
    \benum
        \item {\emphcolor connected} if it cannot be written as the disjoint union of two non-empty open subsets (equivalently it has no non-trivial clopen sets);
        \item {\emphcolor path connected} if every two points in $X$ can be connected by a path (a continuous $\gamma\colon[0,1]\to X$ whose endpoints are the given points);
        \item {\emphcolor locally path connected} if $X$ has a basis of path-connected open subsets.
    \eenum

    Furthermore, $X$'s {\emphcolor connected components} is the set of its maximally connected subsets.
    Its {\emphcolor path components} is the set of its maximally path connected subsets.

\eenum

\bprop

    Let $M$ be a topological manifold, then
    \benum
        \item $M$ is locally path connected,
        \item $M$ is connected iff it is path connected,
        \item The connected and path components of $M$ coincide,
        \item $M$ has countably many components, all of which is an open subset of $M$ and a connected topological manifold.
    \eenum

\eprop

\bproof

    Every point in $M$ is contained in a coordinate ball (a coordinate chart whose image is a ball).
    Since balls are path-connected, the coordinate ball is as well.
    Thus we can form a basis of coordinate balls, which are all path connected.

    $(2)$ and $(3)$ are true in any locally path connected space.
    Furthermore in any locally path connected space, its components are open.
    Thus the components form an open cover of $M$, and must have a countable subcover by second-countability.
    But components are disjoint, and thus form a minimal cover; the cover must be countable already.
    In general an open subset of a manifold is itself a manifold clearly.
    \qed

\eproof

\bdefn

    A topological space $X$ is {\emphcolor locally compact} if every point has a {\emphcolor compact neighborhood}.
    That is to say, for every $x\in X$ there is an open $U$ and compact $K$ such that $x\in U\subseteq K$.

    A {\emphcolor precompact} set in $X$ is a subset $Y\subseteq X$ whose closure is compact.

\edefn

\blemm

    Every topological manifold $M$ has a countable basis of precompact coordinate balls.

\elemm

\bproof

    First consider the case where $M$ can be covered by a single chart, $\phi\colon M\to\hat U\subseteq\bR^n$.
    Consider the set of all balls with rational center and radius $B_q(x)$ such that $B_{q'}(x)\subseteq\hat U$ for some $q'>q$.
    These form a countable precompact basis of $\hat U$ (they are precompact because their closure in $\bR^n$ is contained in the closed ball of the same radius,
    which in turn is contained in $B_{q'}(x)$).
    $\phi$ lifts this precompact basis to a precompact basis of $M$, as desired.

    $M$ can be covered by countably many charts: take all the charts and then take a countable subcover, which exists by second-countability.
    Each of these charts $U_i$ has a countable precompact basis, and if $V\subseteq U_i$ is precompact then it is precompact in $M$.
    Thus the countable union of these countable precompact bases is a countable precompact basis for $M$, as desired.
    \qed

\eproof

\bcoro

    A topological manifold is locally compact.

\ecoro

\bproof

    The basis of precompact coordinate balls gives the desired result: let $p\in U$ be a precompact coordinate ball, then $p\in U\subseteq\bar U$ is compact.
    \qed

\eproof

\bdefn

    Let $X$ be a topological space.
    Then a collection $\c X$ of subsets of $X$ is {\emphcolor locally finite} if every $x\in X$ has a neighborhood that intersects at most finitely many subsets in $\c X$.
    Given a cover $\c U$ of $X$, a {\emphcolor refinement} of $\c U$ is another cover $\c V$ such that every set in $\c V$ is contained in some set in $\c U$.
    $X$ is {\emphcolor paracompact} if every open cover of $X$ admits an open locally finite refinement.

\edefn

If $\c X$ is locally finite, then so is $\cl{\c X}=\set{\cl X}[X\in\c X]$ and we have
$$ \cl{\bigcup\c X} = \bigcup\cl X $$

\bthrm

    Topological manifolds are paracompact.
    Moreso, if $\U$ is an open cover of $M$ and $\B$ is a basis, there is a locally finite open refinement of $\U$ consisting of elements from $\B$.

\ethrm

\bproof

    The above lemma tells us that $M$ can be covered by countably many compact subsets.
    Recall that the union of finitely many compact subsets is compact, and so we can construct an ascending chain of compact subsets $(K_i)_i$ whose union is $M$.
    We can also require that $K_i\subseteq K^\circ_{i+1}$.

    Define $V_i=K_{i+1}-K_i^\circ$ and $W_i=K_{i+2}^\circ-K_{i-1}$ for all $i$ (let $K_i=\emptyset$ for negative $i$).
    Then $V_i$ is compact, and contained in the open set $W_i$.
    For $x\in V_i$ there is some $X_x\in\c X$ containing $x$, and since $\B$ is a basis there is a $B_x\in\B$ such that $x\in B_x\subseteq X_x\cap W_i$.
    Consider then the set of all such $B_x$ as $x$ ranges over $V_i$, this forms an open cover of $V_i$ and thus has a finite subcover.

    The union of all these finite subcovers as $i$ ranges over all integers forms a countable open cover of $M$ that refines $\c X$.
    The finite subcover of $V_i$ consists of sets contained in $W_i$, and $W_i\cap W_{i'}=\emptyset$ except when $\abs{j-j'}\leq2$, thus this is a locally finite open cover
    of $M$ refining $\c X$.
    \qed

\eproof

\subsection{Smooth structures}

Thus far we have restricted our attention to {\it topological\/} manifolds.
But in this course we want to do analysis on manifolds, so we must now consider smooth structures on manifolds.

Recall that a function $f\colon U\to V$ for $U\subseteq\bR^n,V\subseteq\bR^m$ is $C^k$ (for a potentially infinite nonnegative integer $k$) if it is continuous,
and each of its component functions can be continuously differentiated at least $k$ times.
That is to say, $f=(f_1,\dots,f_m)$ and each of its partial derivatives
$$ \frac{\partial^\alpha f}{\partial x_1^{\alpha_1}\cdots\partial x_n^{\alpha_n}} $$
exist and are continious, for $\alpha=\alpha_1+\cdots+\alpha_n\leq k$.

We are mostly concerned with $C^\infty$ functions; those which can be continuously differentiated an arbitrary number of times.
Such functions are also called {\emphcolor smooth}.
A smooth function with a smooth inverse is called a {\emphcolor diffeomorphism}.

\bdefn

    Let $M$ be an $n$-manifold.
    Let $(U,\phi),(V,\psi)$ be two non-disjoint charts.
    Consider then the map $\psi\circ\phi^{-1}\colon\phi(U\cap V)\to\psi(U\cap V)$, this is called the {\emphcolor transition map} from $\phi$ to $\psi$.
    As the composition of homeomorphisms, it itself is a homeomorphism.
    Two charts $(U,\phi),(V,\psi)$ are {\emphcolor smoothly compatible} if they are either disjoint or their transition map is a diffeomorphism.

    An {\emphcolor atlas} for $M$ is a collection of charts whose domains cover $M$.
    A {\emphcolor smooth atlas} for $M$ is an atlas whose charts are all smoothly compatible.

\edefn

To show that an atlas is smooth, we need only show that each transition map is smooth (and not a diffeomorphism).
This is because its inverse is a transition map itself: $(\psi\circ\phi^{-1})^{-1}=\phi\circ\psi^{-1}$.

\bdefn

    Let $M$ be a manifold, then a {\emphcolor smooth structure} on $M$ is a maximal smooth atlas (i.e.\ one not contained in any other atlas).
    A {\emphcolor smooth manifold} is a pair $(M,\A)$ where $\A$ is a smooth structure on a topological manifold $M$.

\edefn

\bprop

    Let $M$ be a topological manifold.
    \benum
        \item An atlas for $M$ is contained in a unique maximal smooth atlas (smooth structure),
        \item Two smooth atlases for $M$ are contained in the same smooth structure iff their union is a smooth atlas.
    \eenum

\eprop

\bproof

    Let $\A$ be a smooth atlas for $M$, and define $\cl\A$ to be the set of all charts smoothly compatible with every chart in $\A$.
    We claim that $\cl\A$ is itself a smooth atlas; i.e.\ that every two charts in it are smoothly compatible.

    Let $(U,\phi),(V,\psi)$ be charts in $\cl\A$, and let $x=\phi(p)\in\phi(U\cap V)$.
    Then take some chart $(W,\theta)$ from $\A$ such that $p\in W$.
    Now $(U,\phi)$ and $(V,\psi)$ are both smoothly compatible with $(W,\theta)$, so $\theta\circ\phi^{-1},\psi\circ\theta^{-1}$ are smooth when defined.
    In particular since $p\in U\cap V\cap W$, $\psi\circ\phi^{-1}$ is smooth, as the composition of these two smooth functions, in a neighborhood of $p$.
    This is true for every $p\in U\cap V$, so $\psi\circ\phi^{-1}$ is smooth as desired.
    Thus $\cl\A$ is indeed a smooth atlas.
    It is also clearly maximal: any chart not in it by definition is not smoothly compatible with some chart in $\A\subseteq\cl\A$.

    If $\B$ is any other smooth structure containing $\A$, then its charts are smoothly compatible with $\A$ and thus by definition $\B\subseteq\cl\A$.
    By maximality, this is equality.
    So we have shown the first point.

    The second point is clear: if $\A,\B$ are contained in the same smooth structure, then their union is as well and thus must be a smooth atlas.
    Conversely if $\A\cup\B$ is a smooth atlas, it is contained in a unique smooth structure, and this smooth structure contains $\A$ and $\B$.
    \qed

\eproof

\bdefn

    Let $\A$ be a smooth atlas for $M$, and $\B$ a smooth atlas for $N$.
    A map $f\colon M\to N$ is {\emphcolor smooth} if for every chart $p\in M$ there are charts $(U,\phi)\in\A,(V,\psi)\in\B$ with $p\in U$ and $f(U)\subseteq V$ such that
    the composite $\psi\circ f\circ\phi^{-1}\colon\phi(U)\to\psi(V)$ is smooth.
    A bijection $f$ is a {\emphcolor diffeomorphism} if it is smooth and has a smooth inverse.

\edefn

Note that this definition should seemingly talk about smooth structures, and not just smooth atlases.
But this is unnecessary:

\bprop

    $f\colon(M,\A)\to(N,\B)$ is smooth iff $f\colon(M,\cl\A)\to(N,\cl\B)$ is.

\eprop

\bproof

    Clearly the right implies the left.
    Now let $(U,\phi)$ be a chart compatible with $\A$ and $(V,\psi)$ compatible with $\B$.
    For $p\in\phi(U)$ let $(W,\theta)$ be a chart in $\A$ with $p$ in its domain, and $(W',\theta')$ a chart in $\B$ with $f(p)$ in its domain.
    Then we can write
    $$ \psi\circ f\circ\phi^{-1} = (\psi\circ{\theta'}^{-1})\circ(\theta'\circ f\circ\theta^{-1})\circ(\theta\circ\phi^{-1}), $$
    a composition of smooth maps, and is thus smooth.
    So $\psi\circ f\circ\phi^{-1}$ is smooth on a neighborhood of any element in its domain, thus smooth as desired.
    \qed

\eproof

Thus we can characterize smoothness by any compatible smooth atlas.
An immediate result is the following (expanding definitions):

\blemm

    Two smooth atlases $\A,\B$ for $M$ are smoothly compatible iff the identity $(M,\A)\to(M,\B)$ is a diffeomorphism.

\elemm

Note that two atlases may give diffeomorphic manifolds without them being smoothly compatible.
Necessarily the diffeomorphism must be different from the identity.

So instead of talking about maximal smooth atlases, we can concentrate on ordinary smooth atlases, while keeping in mind that distinct smooth atlases can give the same smooth structure.

\bexam

    Clearly Euclidean space $\bR^n$ can be given a natural smooth structure: the one induced by the atlas $\id_{\bR^n}$.
    This is called the {\emphcolor standard smooth structure}.
    These are not the only smooth structures possible: $\phi(x)=x^3$ is a homeomorphism $\bR\to\bR$ and thus defines a smooth atlas
    $(\bR,\phi)$.
    Since $\phi$ is not a diffeomorphism, it is not smoothly compatible with the standard structure.
    Despite this $(\bR,\phi)$ is diffeomorphic to the standard structure, just via a non-identity map.

\eexam

\bexam

    Let $V$ be a finite-dimensional real vector space, so that it is isomorphic to $\bR^n$.
    Any two norms on $V$ are equivalent, i.e.\ they induce the same topology.
    Thus $V$ can be given a smooth structure by lifting $\bR^n$'s, and this agrees with any norm-induced topology.

    For example, the space of matrices $M_{n\times m}(\bR)$ has a natural smooth structure.
    Similarly for $M_{n\times m}(k)$ for any finite-dimensional $\bR$-algebra $k$ (e.g.\ the complex numbers or quaternions).

\eexam

\bexam

    Let $M$ be a manifold and $U\subseteq M$ be open.
    Then for an atlas $\A$ of $M$, we can restrict it to $U$ by
    $$ \A_U = \set{(U\cap V,\phi|_{U\cap V})}[(V,\phi)\in\A] . $$
    If $\A$ is a smooth structure, this is just all of the charts $(V,\phi)\in\A$ such that $V\subseteq U$.
    This is called an {\emphcolor open submanifold} of $M$.

\eexam

\bexam

    Consider the general linear group $\gl_n(\bR)\subseteq M_n(\bR)$.
    This is the preimage of $\bR-0$ under $\det$, and is thus an open subset of $M_n(\bR)$.
    As discussed, this means that $\gl_n(\bR)$ is an open submanifold of $M_n(\bR)$.

    More generally, consider the subspace of $M_{n\times m}(\bR)$ of matrices of full rank.
    This is also an open subspace: if $A$ has full rank, then consider its submatrix given by its linearly independent columns (or rows).
    We can then define a continuous map $\det\colon M_{n\times m}(\bR)\to\bR$ by mapping a matrix to the determinant of this square matrix.
    Then the inverse image of $\bR-0$ is an open neighborhood of $A$, and all matrices in it are of full rank.
    So the set of matrices of full rank is indeed an open submanifold.

\eexam

\bexam

    We showed previously that $S^n$ is a topological $n$-manifold with charts $(U_i^\pm,\phi_i^\pm)$ for $i=1,\dots,n+1$ where
    $\phi_i^\pm$ removes the $i$th coordinate, and its inverse maps $\bar x$ to the coordinate with $\sqrt{1-\norm x^2}$.
    Thus for $i\neq j$, the transition $\phi_i^\pm\circ(\phi_j^\pm)^{-1}$ maps $x$ to the coordinate with $x_i$ removed and
    $\pm\sqrt{1-\norm x^2}$ in the $j$th coordinate.
    This is smooth.
    For $i=j$ we just get the identity, which is smooth.

    Thus $S^n$ is a smooth $n$-manifold.

\eexam

\bexam

    Similar to before the $n$-dimensional topological manifold $\RP^n$ is smooth.
    Its atlas is given by maps $(U_i,\phi_i)$ where $U_i\subseteq\RP^n$ is the set of all $[x]$ where $x_i\neq0$ and
    $\phi_i[x]$ is $1/x^i$ times the vector $x$ with its $i$th coordinate removed.
    Then $\phi_j\circ\phi_i^{-1}$ maps $x$ to $1/x^j$ times the vector obtained by removing the $j$th coordinate of $x$ and adding
    $1$ to its $i$th coordinate.
    This is smooth.

\eexam

Recall that if $M_1,\dots,M_n$ were topological manifolds with charts $(U_i,\phi_i)$, then $M_1\times\cdots\times M_n$ was a topological
manifold with charts $(U_1\times\cdots\times U_n,\phi_1\times\cdots\phi_n)$.
If these are smooth atlases, then these charts are smoothly compatible, since
$$ (\psi_1\times\cdots\times\psi_n)\circ(\phi_1\times\cdots\times\phi_n)^{-1} = (\psi_1\circ\phi_1^{-1})\times\cdots\times(\psi_n\circ\phi_n^{-1}) $$
which is smooth.
So we have proven the following.

\bprop

    The finite product of smooth manifolds is a smooth manifold, whose dimension is the sum of the dimensions of the manifolds.

\eprop

For example, the $n$-torus is a smooth manifold of dimension $n$.

\blemm

    Let $M$ be a set and suppose we have a collection of subsets $\set{U_i}_{i\in I}$ with maps $\phi_i\colon U_i\to\bR^n$
    for some $n$.
    Further suppose the following:
    \benum
        \item for each $i$, $\phi_i(U_i)$ is open and $\phi_i$ is injective;
        \item for each $i,j$, $\phi_i(U_i\cap U_j)$ and $\phi_j(U_i\cap U_j)$ are open in $\bR^n$;
        \item for $\phi_i\cap\phi_j^{-1}\colon\phi_i(U_i\cap U_j)\to\phi_j(U_i\cap U_j)$ is smooth;
        \item $\set{U_i}_i$ has a countable subcover;
        \item whenever $p\neq q\in M$ are distinct, there is either a $U_i$ containing them both, or disjoint $U_i,U_j$ such that
        $p\in U_i$ and $q\in U_j$,
    \eenum
    then $M$ has a unique smooth manifold structure such that $(U_i,\phi_i)$ are all smooth charts.

\elemm

\bproof

    We first need to give $M$ a topology: take $\phi_i^{-1}(V)$ to be a basis for $V\subseteq\bR^n$ open.
    In particular $U_i=\phi_i^{-1}(\bR^n)$ is open.
    To show that this is a basis, we need to show that the intersection $\phi_i^{-1}(V_i)\cap\phi_j^{-1}(V_j)$ can be written as a union of
    base sets:
    $$ \phi_i^{-1}(V_i) \cap \phi_j^{-1}(V_j) = \phi_i^{-1}(V_i\cap(\phi_j\cap\phi_i)^{-1}(V_j)) , $$
    which is a base set, since $(\phi_j\cap\phi_i)^{-1}(V_j)$ is an open subset of the open subset $\phi_i(U_i\cap U_j)$ and is thus open.
    Note that this is the unique topology making each $\phi_i$ a homeomorphism; thus far we have made no choices.

    Then $M$ is Hausdorff: for $p\neq q$ either they are in the same $U_i$ which is Hausdorff (as it is Euclidean), or in distinct $U_i,U_j$
    which are open.
    And $M$ is second countable: as $U_i$ are Euclidean, they are second countable, and $M$ is thus second countable as the union of countably
    many $U_i$s.

    So we have shown that $M$ is a topological $n$-manifold.
    And then $(U_i,\phi_i)$ are all smoothly compatible charts because we assumed the transition maps were smooth.
    \qed

\eproof

\bnote

    We could state a similar result for topological manifolds, whose proof is the same as above.
    Instead of assuming the transition maps are smooth, we relax this condition to just require they are continuous.
    Then $M$ has a unique topological manifold structure making $(U_i,\phi_i)$ all charts.

\enote

\bexam

    Let $K$ be a finite-dimensional real vector space of dimension $n$.
    For $0\leq k\leq n$, define $G_k(V)$ to be the set of all $k$-dimensional subspaces of $V$.
    Using the abovbe lemma, we can turn this into a smooth manifold called the {\emphcolor Grassmanian manifold}.
    For $U\leq V$ of dimension $k$, let $W$ be a complementary subspace such that $V=U\oplus W$.
    Then define $\Gamma_{U,W}\colon\hom(U,W)\to G_k(V)$ by $\Gamma_{U,W}(L)=\set{u+Lu}[u\in U]$.

    Then $(\Gamma_{U,W}\hom(U,W),\Gamma_{U,W}^{-1})$ satisfy the conditions above, turning $G_k(V)$ into a $k(n-k)$-dimensional smooth
    manifold (since $\hom(U,W)$ has dimension $k(n-k)$).
    The proof of this is not trivial, and can be found in my homework solutions.

    For $k=1$, notice that $G_1(\bR^{n+1})\cong\RP^n$.

\eexam

\subsection{Manifolds with boundary}

Define the {\emphcolor closed $n$-dimensional upper half-space} $\bH^n\subseteq\bR^n$ as the set of all points whose $n$th coordinate is
nonnegative:
$$ \bH^n = \set{\bar x\in\bR^n}[x_n\geq0] . $$
Its interior and boundary are
$$ \Int\bH^n = \set{\bar x\in\bR^n}[x_n>0],\qquad \partial\bH^n = \set{\bar x\in\bR^n}[x_n=0] . $$

\bdefn

    An $n$-dimensional {\emphcolor (topological) manifold with boundary} is a second-countable Hausdorff space $M$ such that every
    point in $M$ lies in a neighborhood homeomorphic to an open subset of $\bH^n$.

    A {\emphcolor chart} for $M$ is a pair $(U,\phi)$ for $U\subseteq M$ open and $\phi\colon U\to\phi(U)$ a homeomorphism to an open subset $\phi(U)$
    of $\bH^n$.
    Call a chart a {\emphcolor boundary chart} if $\phi(U)\cap\partial\bH^n\neq\emptyset$, and an {\emphcolor interior chart} otherwise
    (if its image lies in the interior of $\bH^n$).
    An {\emphcolor atlas} for $M$ is a collection of charts whose domains cover $M$.

    Call $p\in M$ a {\emphcolor boundary point} if it lies in the domain of some boundary chart $(U,\phi)$ such that $\phi(p)\in\partial\bH^n$, and
    an {\emphcolor interior point} if it lies in the domain of some interior chart.
    Define the set of boundary and interior points to be
    $$ \partial M\hbox{ and }\Int M $$
    respectively.

\edefn

Open subsets of $\bH^n$ are of the form $\bH^n\cap U$ for an open subset $U\subseteq\bR^n$.
In particular we can construct a basis of $\bH^n$ by $\bH^n\cap B_r(x)$.
We can refine this basis by considering only balls of the form $B_r(x)$ for $B_r(x)\subseteq\bH^n$ and $B_r(x)\cap\bH^n$ for $x\in\partial\bH^n$,
called {\emphcolor half-balls}.
If $(U,\phi)$ is a chart with $\phi(U)$ a half-ball, call $U$ a {\emphcolor coordinate half-ball}.

Note that every point is either a boundary point or an interior point: suppose $p\in M$ is neither on the boundary or interior.
So it must be in the domain of a boundary chart $(U,\phi)$, but $\phi(p)\notin\partial\bH^n$.
But then we can take a ball around $\phi(p)$ so that it doesn't intersect the boundary of $\bH^n$, and lift it to $M$; this gives an interior
chart with $p$ in its domain, so $p$ is an interior point.

Note that an ordinary manifold is also a manifold with boundary, but without boundary points.
This is because open subsets in $\bR^n$ are all homeomorphic to open subsets of $\bH^n$: for example a ball $B_r(x)\subseteq\bR^n$ is homeomorphic
to the ball $B_r(x')\subseteq\Int\bH^n$ where $x'$ is just $x$ but the last coordinate is raised so that $B_r(x')$ lies within the interior of the
half-space.
Thus we can view manifolds with boundaries as a generalization of ordinary manifolds.
We will still say ``manifolds with or without boundary'' when discussing manifolds with boundary to emphasize the point that a result holds
for ordinary manifolds as well.

As it turns out, the boundary and interior of a manifold are topological invariants.

\bthrm

    If $M$ is a topological manifold with boundary, then every point is an interior or boundary point, but not poth.
    That is,
    $$ M = \partial M \sqcup \Int M . $$

\ethrm

This can be proven using standard tools from algebraic topology.

Note that the concept of a manifold boundary and interior is distinct (yet related) to the notion of a topological boundary and interior.
The notion of a topological boundary and interior is {\it relative\/}: it is defined relative to a larger topological space.
While manifold boundaries and interior are not.
To see how this can manifest, consider the manifold $\bR$.
It has an empty boundary as a manifold, and also as a topological space (as is true for all topological spaces, relative to themselves).
But viewed as a topological subspace of $\bR^2$ say, its boundary is itself.

Thus it is important to distinguish manifold and topological boundaries and interiors.

\bprop

    Let $M$ be a topological $n$-manifold with boundary.
    \benum
        \item $\Int M$ is an open subset of $M$ and a topological $n$-manifold, without boundary.
        \item $\partial M$ is a closed subset of $M$, and a topological $(n-1)$-manifold, without boundary.
        \item $M$ is a topological manifold iff $\partial M$ is empty.
    \eenum

\eprop

\bproof

    For $p\in\Int M$, the domain of its interior chart also lies in $\Int M$ by definition, so $\Int M$ must be open in $M$.
    It is clearly a topological $n$-manifold, as its charts are to $\bR^n$.

    Since $\partial M$ is the complement of the open $\Int M$, it is closed.
    For $p\in\partial M$ consider its boundary chart mapping $p$ to $\partial\bH^n$, $(U,\phi)$.
    Then the restriction of $\phi$ to $U\cap\partial M$ is a homeomorphism to an open subset of $\partial\bH^n\cong\bR^{n-1}$.
    Indeed, $\phi$ must map boundary points to boundary points, otherwise we could lift an open ball to show that a boundary point is also
    an interior point.
    So $\partial M$ is locally Euclidean of dimension $n-1$, as desired.

    Point (3) follows directly from the invariance of the boundary and point (1).
    \qed

\eproof

$\Int M$ being a submanifold of $M$ is true for all open subsets.

\bprop

    Let $M$ be a topological $n$-manifold with boundary, and $U\subseteq M$ open.
    Then $U$ is a topological $n$-manifold with boundary $\partial U=U\cap\partial M$.

\eprop

\bproof

    For $p\in U$ consider its chart in $M$ and restrict it to $U$.
    $p\in U$ is then a boundary point in $U$ iff it is a boundary point in $M$.
    \qed

\eproof

\bprop

    Let $M$ be a topological $n$-manifold with boundary.
    Then
    \benum
        \item $M$ has a countable basis of precompact coordinate balls and half-balls.
        \item $M$ is locally compact.
        \item $M$ is paracompact.
        \item $M$ is locally path-connected.
        \item $M$ has countably many components, each of which is an open subset of $M$
    \eenum

\eprop

Recall that a function $f\colon A\to\bR^m$ for $A\subseteq\bR^n$ arbitrary is {\emphcolor smooth} if $A$ has an open neighborhood $U$ such that
$f$ can be extended to $U\to\bR^m$ and this is smooth.
In particular if $U\subseteq\bH^n$ is open, then $f\colon U\to\bR^m$ is smooth if there is an open subset $\tilde U\subseteq\bR^n$ such that
$U=\tilde U\cap\bH^n$ and a lift $\tilde f\colon W\to\bR^m$ of $f$, which is smooth.
Notice that if $f$ is smooth, then $f$'s restriction to $U\cap\Int M$ is smooth in the usual sense.

So say we have a topological manifold with boundary $M$, and let $(U,\phi),(V,\psi)$ be charts.
Then we have the transition map $\psi\circ\phi^{-1}\colon\phi(U\cap V)\to\psi(U\cap V)$, which is a homeomorphism.
To say that this is smooth is to say $\phi(U\cap V)$ has an open neighborhood $W\subseteq\bR^n$ to which we can extend $\psi\circ\phi^{-1}$,
and this lift is smooth.

\bdefn

    Let $M$ be a manifold with boundary.
    Then an {\emphcolor atlas} for $M$ is a set of charts $(U_i,\phi_i)$ whose domain covers $M$.
    A {\emphcolor smooth atlas} for $M$ is an atlas for which all the transition maps $\phi_j\circ\phi_i^{-1}$ are smooth in the sense explained
    above.
    A {\emphcolor smooth structure} for $M$ is a maximally smooth atlas.
    A {\emphcolor smooth manifold with boundary} is a manifold $M$ along with a choice of a smooth structure $\A$.

\edefn

As before, a smooth atlas uniquely determines a unique smooth structure.

\bdefn

    Let $M,N$ be manifolds with boundary.
    A map $F\colon M\to N$ is {\emphcolor smooth} if for every point $p\in M$, there are smooth charts $(U,\phi)\in M$ and $(V,\psi)\in N$
    such that $p\in U$ and $f(U)\subseteq V$ and $\psi\circ f\circ\phi^{-1}\circ\phi(U)\to\psi(V)$ is smooth.

    A {\emphcolor diffeomorphism} is a smooth map with a smooth inverse.

\edefn

As before, a function is smooth relative to a smooth atlas iff it is smooth relative to the smooth structure induced by that atlas.
And two smooth atlases define the same smooth structure iff the identity is a diffeomorphism with respect to them.

\bprop

    If $M$ is a smooth manifold with boundary, and $U\subseteq M$ is open, then it is also a smooth manifold with boundary.

\eprop

\bproof

    Consider the construction above showing $U$ is a topological manifold with boundary, the smoothness of the chart given is evident.
    \qed

\eproof

\bnote

    Consider our equivalent definition of a smooth manifold via a covering of a set by injections into $\bR^n$.
    If we replace $\bR^n$ with $\bH^n$, then this gives an equivalent definition of smooth manifolds with boundary.

\enote

Interestingly, the product of manifolds with boundary is not necessarily a manifold with boundary.
This is because the product of half-spaces is not a half-space.
But we can take the product of a manifold with boundary with a manifold (without boundary).

\bprop

    Let $M_1,\dots,M_n$ be manifolds (without boundary), and $N$ a manifold with boundary.
    Then their direct product $M_1\times\cdots\times M_n\times N$ is a manifold with boundary
    $$ \partial(M_1\times\cdots\times M_n\times N) = M_1\times\cdots\times M_n\times\partial N . $$

\eprop

